\documentclass[11pt]{article}
\usepackage[margin=1.05in]{geometry}
\usepackage{amsmath,amssymb}
\usepackage{booktabs}
\usepackage{enumitem}
\usepackage[hidelinks]{hyperref}

\title{Segment-Dependent Chunk Length: Stages 1, 2, and 3\\
\large One Combined Report}
\author{Aryan Goyal}
\date{July 23, 2026}

\begin{document}
\maketitle

\begin{abstract}
\noindent
An action-chunking policy predicts several future actions at once and then
decides how many of them to execute before looking again. That number is
usually a fixed constant. This project asks whether it should instead depend
on the moment: commit a long chunk where the physics absorbs errors, and
replan sooner where it amplifies them. This report combines the three stages
of that investigation. Stage 1 tests the idea on short benchmark tasks, first
with a perfect oracle signal and then with a learned predictor. Stage 2 tests
it on long-horizon tasks that alternate between the two regimes, which is
where the theory says the win should appear. Stage 3 proposes a constructed
task where the regime mix becomes an experimental knob. Stage 1 is complete,
Stage 2 is in its final evaluation, and Stage 3 is a design.
\end{abstract}

\tableofcontents
\newpage

%=====================================================================
\section{The hypothesis and the stage map}

Executing $k$ actions without looking has two competing costs. If $k$ is too
long, the robot is blind through a moment where a small error grows into a
failure. If $k$ is too short, the policy queries itself constantly, and every
query injects a fresh sample of the policy's own error. The right $k$
therefore depends on whether the current moment amplifies errors or absorbs
them. We measure that property directly in simulation by perturbing one
action and watching whether the trajectory deviation grows (open-loop
unstable) or not (open-loop stable). The labeling procedure is documented
separately in the labeling report.

Each stage answers one question.

\begin{center}
\small
\begin{tabular}{@{}llll@{}}
\toprule
Stage & Question & Signal source & Status \\
\midrule
1a & Does chunk length matter, and in which direction? & perfect oracle & done \\
1b & Can a model read the regime from camera frames? & trained predictor & done \\
1c & Does the predicted signal help in rollouts? & trained predictor & done \\
2  & Does it help on long alternating tasks? & predictor, zero-shot & running \\
3  & Does the gain grow with more alternations? & task construction & proposed \\
\bottomrule
\end{tabular}
\end{center}

The logic is sequential. If a perfect signal does not help (1a), no predictor
can. If the regime cannot be read from frames (1b), the rule cannot leave the
simulator. If the predicted signal helps in rollouts (1c), the remaining
question is whether the benchmark tasks were ever the right test, which is
Stage 2. Stage 3 exists because Stage 2 cannot control how much of each
episode is stable transport and how much is unstable contact, and that mix is
what decides how much room the switch has to win.

%=====================================================================
\section{Stage 1a: does chunk length matter?}

\subsection{Setup}

Two image-based diffusion policies trained on robomimic proficient-human
data: square (epoch 950, success 0.90) and tool\_hang (epoch 1200, success
0.90). Nothing about the policy changes between cells. Only the execution
rule changes.

\begin{itemize}[leftmargin=1.4em]
\item \textbf{Fixed cells.} One chunk length for the whole episode,
$k \in \{1, 2, 4, 8, 16\}$. Note that $k{=}1$ is ordinary per-step control.
\item \textbf{Oracle cells.} Two chunk lengths
$(k_{\mathrm{free}}, k_{\mathrm{contact}})$, switched by reading the
simulator's contact table at every query, with a two-query hysteresis. Grid:
$(1,4)$, $(1,16)$, $(4,1)$, $(8,1)$, $(8,4)$, $(16,1)$, $(16,4)$.
\end{itemize}

Contact is the oracle signal rather than the labels themselves because the
true labels exist only on demonstration states, and a rollout leaves those
states immediately. Contact is the best label-aligned observable available:
the labeling report shows fixture contact above gripper contact above free
space in mean divergence on every dataset. Each cell is 50 episodes, repeated
at three environment seeds, so 150 episodes per cell and a standard error near
0.033.

\subsection{Results}

\begin{center}
\small
\begin{tabular}{@{}lcc@{}}
\toprule
Cell & Square & Tool\_hang \\
\midrule
fixed $k{=}1$  & 0.793 & 0.667 \\
fixed $k{=}2$  & 0.820 & 0.740 \\
fixed $k{=}4$  & 0.807 & 0.793 \\
fixed $k{=}8$  & 0.833 & 0.800 \\
fixed $k{=}16$ & 0.827 & \textbf{0.807} \\
\midrule
oracle $(8,4)$  & \textbf{0.873} & \textbf{0.807} \\
oracle $(16,4)$ & 0.793 & 0.780 \\
oracle $(1,4)$  & 0.773 & 0.780 \\
oracle $(4,1)$  & 0.793 & 0.600 \\
oracle $(8,1)$  & 0.820 & 0.667 \\
oracle $(16,1)$ & 0.780 & 0.633 \\
\bottomrule
\end{tabular}
\end{center}

Three findings, all on tool\_hang, the contact-heavy task where the labels
say the effect should appear.

\begin{enumerate}[leftmargin=1.6em]
\item \textbf{Short chunks hurt.} Fixed $k{=}1$ reaches 0.667 against 0.807
for $k{=}16$, a gap of 2.8 standard errors. Chunk lengths 4 through 16 form a
plateau; the damage is concentrated at $k \leq 2$.
\item \textbf{Replanning every step during contact is the worst rule
measured.} The three oracle cells that replan per step in contact reach
0.600, 0.633, and 0.667, which is 3 to 5 sigma below the plateau, in three
independent cells. This is the direction the closed-loop labels predicted:
the policy's own corrections inject error, and contact amplifies it.
\item \textbf{The binary oracle never beats the best fixed $k$.} The best
oracle cell $(8,4)$ reaches 0.807 and exactly ties fixed $k{=}16$.
\end{enumerate}

Square stays flat across the whole grid (0.773 to 0.873), which is the
expected null for a task whose policy is strong at every $k$.

\subsection{Why the oracle only tied}

The switching signal carried almost no information. Diagnostics show the
oracle judged 84 to 88 percent of its tool\_hang decisions to be in contact,
so a binary switch is effectively constant and collapses to fixed $k$ at its
contact value. This also explains the cell-by-cell pattern: every cell with
contact value 1 lands near fixed $k{=}1$, and every cell with contact value 4
lands near fixed $k{=}4$. What the labels actually contain is a continuous
grade of instability, and binarizing it threw that away. Recovering it is the
purpose of the predictor.

A secondary observation supports the injection mechanism: mean steps among
successful episodes falls monotonically with commitment, from 167 to 145 on
square and 480 to 436 on tool\_hang. Per-step replanning dithers, commitment
produces straighter trajectories.

%=====================================================================
\section{Stage 1b: can the regime be read from frames?}

\subsection{Setup}

A frozen V-JEPA 2 video encoder (ViT-L, 300M parameters, never updated) plus
a small trained head (attentive pooling, proprioception fusion, two-layer
MLP). Input per stamp is the last 16 camera frames and a 9-dimensional
proprioception vector. Output is the probability the moment is unstable and a
predicted divergence rate. Features are cached once, so head training is
minutes per epoch. Training pool: about 155{,}000 labeled stamps across the
four robomimic tasks, two MimicGen sets, ten LIBERO-Long files, and policy
rollout states. Splits are by demonstration, never by stamp, because stamps 2
steps apart are nearly identical.

\subsection{Results}

\begin{center}
\small
\begin{tabular}{@{}lccl@{}}
\toprule
Run & AUROC all & AUROC confident & Note \\
\midrule
1: pooled head            & 0.699 & 0.868 & the head used downstream \\
2: square only            & 0.768 & 0.719 & single-task training \\
2: square to tool\_hang   & 0.562 & 0.465 & transfer test \\
3b: pooled plus LIBERO    & 0.720 & 0.893 & pinned validation \\
4: DINOv2, single frame   & 0.703 & 0.859 & motion ablation \\
5: closed-loop labels     & 0.627 & n/a   & diagnostic head \\
\bottomrule
\end{tabular}
\end{center}

``Confident'' restricts scoring to stamps outside the deadband, which are the
only stamps whose binary label is proven rather than ambiguous. Four
conclusions.

\begin{enumerate}[leftmargin=1.6em]
\item \textbf{The pooled head works where the labels are confident.} 0.868
means the head separates proven-unstable from proven-stable moments from
frames plus proprioception. The lower all-stamps number is expected, since
deadband stamps are ambiguous in the ground truth itself.
\item \textbf{Single-task training does not transfer.} Square to tool\_hang
lands at chance or below. A head trained on one task learns task-specific
visual cues, not stability. Every task family we care about must contribute
labeled data to the pool.
\item \textbf{Motion context adds almost nothing.} A single-frame DINOv2 head
matches the 16-frame video head to within 0.009. The signal the head reads is
scene configuration at the query moment, not motion. A cheaper single-frame
encoder would serve the switch nearly as well.
\item \textbf{Closed-loop expansion is only weakly predictable.} 0.627, well
above chance but nowhere near deployable. This is consistent with the
mechanism: closed-loop expansion depends on how the policy reacts, which is
not fully visible in the scene. Nothing downstream uses this head.
\end{enumerate}

One methodological note worth carrying forward. The first LIBERO comparison
drew its train and validation split randomly over the combined pool, which
leaked most of the baseline's validation demonstrations into training and
showed a fake improvement of 0.968. Pinning the validation demonstrations
showed the true gain, 0.018. Every pool-variant comparison now pins the
split.

%=====================================================================
\section{Stage 1c: does the predicted signal help in rollouts?}

The frozen pooled head goes into the control loop. At every replan the last
16 frames and current proprioception pass through the encoder and head, and
the predicted divergence rate picks the chunk length: 16 if at or below the
training threshold 0.0355, and 4 if above it. Inference runs in a separate
process, about 100 ms per replan. Three seeds pooled, 150 episodes per cell.

\begin{center}
\small
\begin{tabular}{@{}lcccc@{}}
\toprule
 & Predictor $(16,4)$ & Best fixed $k$ & Best oracle cell & Fixed $k{=}1$ \\
\midrule
square     & 0.767 & 0.833 & 0.873 & 0.793 \\
tool\_hang & \textbf{0.800} & 0.807 & 0.807 & 0.667 \\
\bottomrule
\end{tabular}
\end{center}

On tool\_hang the predictor statistically ties both the best fixed $k$ and the
best oracle cell, while avoiding the 0.60 to 0.67 trap of replanning through
contact. Its diagnostics show it is not merely reproducing the contact
signal: it calls unstable on 71.6 percent of decisions against the oracle's
84 to 88 percent contact rate, and its mean executed chunk is 11.4. It
discriminates within contact, committing through contact stretches it judges
benign, which is the continuous-signal advantage Stage 1a said a binary
oracle could not express. On square the predictor sits at the bottom of the
flat band, which is the expected null on a ceiling task.

What Stage 1 did not show is the switch beating the best fixed $k$. Both
tasks are short and 65 to 88 percent contact, so a single fixed $k$ is
already near optimal and a switch has little room. A tie is close to the
ceiling available on these tasks. Testing the idea properly requires tasks
that spend substantial time in both regimes, which is Stage 2.

%=====================================================================
\section{Stage 2: long-horizon alternating tasks}

\subsection{Setup}

Four MimicGen tasks whose episodes are long and alternate between free-space
transport and contact-rich stages: three\_piece\_assembly, nut\_assembly,
kitchen, and coffee\_preparation, all D0, 1{,}000 demonstrations each, mean
episode lengths 337 to 689 steps. Each trains an image-based diffusion policy
for 2{,}000 epochs, which is 200{,}000 gradient steps.

One engineering constraint shapes everything. MimicGen environments require
robosuite 1.4, and our diffusion policy stack requires robosuite 1.5, and the
two cannot coexist in one process. The solution is a strict split: dataset
conversion and environment ownership in one virtual environment, policy
training and inference in another, with a bridge process that streams
observations and actions between them over a pipe. All evaluation is
therefore post hoc rather than during training.

The predictor head is applied zero-shot, with no relabeling of these tasks. A
distribution check on MimicGen frames confirmed the head is not degenerate
there: predicted divergence spans the threshold on all four tasks, with 11 to
55 percent of stamps above it, and the ordering is plausible, since
coffee\_preparation reads most unstable and kitchen most stable. Relabeling
remains the fallback if the zero-shot switch underperforms.

\subsection{Checkpoint selection}

All four trainings finished at epoch 2{,}000. Selection evaluates five
checkpoints per task at 25 episodes each, seed 0, with horizons set per task
from the demonstration lengths (560, 650, 980, and 1{,}140 steps). The rule
is highest success rate, ties broken toward the later epoch.

\begin{center}
\small
\begin{tabular}{@{}lccccc@{}}
\toprule
 & ep.\ 600 & ep.\ 1000 & ep.\ 1400 & ep.\ 1700 & ep.\ 2000 \\
\midrule
kitchen                & 0.96 & 1.00 & 1.00 & 1.00 & \textbf{1.00} \\
three\_piece\_assembly & 0.48 & 0.52 & 0.64 & \textbf{0.80} & 0.68 \\
nut\_assembly          & 0.12 & 0.00 & 0.32 & 0.32 & \textbf{0.36} \\
coffee\_preparation    & 0.04 & running & running & running & pending \\
\bottomrule
\end{tabular}
\end{center}

Kitchen is at ceiling, three\_piece\_assembly is strong, and nut\_assembly is
weak but usable. Coffee\_preparation is the weakest: partial counts from the
running cells put it near 0.15 to 0.25 rather than climbing, so its cells will
carry wide error bars and it may not carry a comparison on its own. The
four-task design anticipated one weak task.

\subsection{The experiment, pending}

On the selected checkpoints, fixed $k \in \{1, 4, 8, 16\}$ plus the predictor
switch $(16,4)$, 50 episodes per cell at seed 0, same policy weights across
all cells. Twenty cells across eight GPUs, roughly a day of compute. The
runner is built and smoke-tested end to end, including the three-process path
through environment, policy, and predictor. It waits on the last selection
cells.

The prediction under test is specific: on these tasks the switch should beat
every fixed $k$, because alternating structure forces any fixed choice to pay
one of the two costs at some point in every episode.

\subsection{Incidents worth recording}

Two training runs were lost and rebuilt. The first loss looked like a random
disk event; the second exposed the cause. This virtual machine reprovisions
its scratch disk on every machine stop, so both losses were the machine being
stopped, not a fault. A separate loss came from a relaunch script that
answered an overwrite prompt automatically and deleted the surviving
checkpoint grids. The process changes are in place: the launcher refuses to
start when an experiment directory exists, checkpoint grids are copied before
any evaluation touches them, and power-off is left to a watchdog that fires
only when trainings, selection, and the experiment have all reported
completion. Net cost was about two days.

%=====================================================================
\section{Stage 3: a constructed regime-switching task}

Stage 2 tests the switch on natural long-horizon tasks, but it cannot control
how much of each episode is stable transport and how much is unstable
contact, and that mix decides how much room the switch has. Stage 3 removes
the limitation by constructing the task: a repeated transport-and-insert task
in robosuite with $N$ pegs spaced far apart, so each cycle is a long
free-space transport followed by a tight insertion, repeated $N$ times. The
two quantities the theory cares about become knobs, namely the number of
alternations $N$ (swept over 1, 2, 4) and the ratio of stable to unstable
segment length (set by spacing and tolerance). The prediction is that the gap
between the switched policy and the best fixed $k$ grows with $N$, because
every added alternation forces a fixed $k$ into one more compromise. The
comparison is the same policy at its best fixed $k$, found by a full sweep,
against oracle switching and then predictor switching. Because we define the
task, regime boundaries are known by construction, which gives a free and
perfect oracle, a ground-truth check of the perturbation labeler against the
designed boundaries, and clean predictor training labels. Demonstrations come
from MimicGen: author the task, record about ten source demonstrations, and
generate one thousand. The known weakness is that a bespoke task invites the
criticism that it was engineered to win, so Stage 3 is positioned as the
mechanism experiment that explains and quantifies the effect while Stage 2
remains the naturalistic evidence. It starts after the Stage 2 results are
in, since those results decide whether it quantifies how the win scales or
diagnoses which regime mix the hypothesis needs.

%=====================================================================
\section{Where things stand}

What is established. Chunk length matters, and it matters in the direction
opposite to the naive control-theory prescription: replanning fastest where
the system is least stable is the worst rule measured, at 3 to 5 sigma across
independent cells. The mechanism is confirmed at two levels, since the
closed-loop labels show the policy's own corrections inject error at 90 to 99
percent of measured states, and the rollouts show per-step replanning losing
0.14 to 0.21 success on the contact-heavy task. The regime is readable from
vision at AUROC 0.868 on confident stamps, and the resulting switch runs in a
real control loop at about 100 ms per decision.

What is not established. The switch has not yet beaten the best fixed chunk
length. On the Stage 1 tasks it ties, and the diagnostics explain why: those
tasks are short and mostly contact, so a fixed $k$ is already near optimal
and a tie is close to the available ceiling. That is consistent with the
hypothesis but is not evidence for it.

What decides it. The Stage 2 experiment, on tasks chosen so that no fixed $k$
can be right throughout. Its preconditions are now met, since three of four
base policies are strong enough to make the comparison meaningful and the
predictor is not degenerate on these scenes. If the switch wins there, the
claim is demonstrated on naturalistic long-horizon tasks and Stage 3
quantifies how the win scales with the number of alternations. If it ties
again, that is a genuine negative result, and Stage 3 becomes the diagnostic
that isolates which regime mix the hypothesis actually needs.

\end{document}
